\par\Needspace{6\baselineskip}\medskip
\begin{center}\hypertarget{native-200-article150}{}\pdfbookmark[2]{150.}{article-150}150.\end{center}
We now illustrate this method with several examples worked out in detail. They will also show very clearly how widely it applies and how readily and quickly it always leads to the desired result\footnote{It is a misuse of words to call one method \textit{more or less exact} than another. A method can count as solving the problem only if it at least makes it possible to attain any degree of precision. Hence one method surpasses another only insofar as the \textit{same} degree of precision can be reached more quickly and with less labor by one, and more slowly and with more labor by the other.}.

The new planet Juno will furnish our \textit{first} example. For this purpose we choose the following observations, made at Greenwich and communicated to us by the distinguished \textsc{Maskelyne}.
\SourceLeaf{201}
\begin{center}
\small\setlength{\tabcolsep}{5pt}\renewcommand{\arraystretch}{1.2}
\begin{tabular}{|r|c|c|}
\hline
Greenwich mean time & App. right ascension & App. south decl.\\
\hline
1804 Oct. 5.\ $10^{\mathrm{h}}\,51^{\mathrm{m}}\,6^{\mathrm{s}}$ & $357^{0}\,10'\,22''{,}35$ & $6^{0}\,40'\,8''$\\
17.\quad $9\quad58\quad10$ & $355\quad43\quad45{,}30$ & $8\quad47\quad25$\\
27.\quad $9\quad16\quad41$ & $355\quad11\quad10{,}95$ & $10\quad2\quad28$\\
\hline
\end{tabular}
\end{center}

From the solar tables we obtain the following for the same times
\begin{center}
\footnotesize\setlength{\tabcolsep}{3pt}\renewcommand{\arraystretch}{1.2}
\begin{tabular}{|r|c|c|c|c|c|}
\hline
 & \shortstack{Solar long. from\\app. equinox} & Nutation & \shortstack{Distance from\\the Earth} & \shortstack{Solar\\latitude} & \shortstack{App. obliquity\\of the ecliptic}\\
\hline
Oct. 5 & $192^{0}\,28'\,53''{,}71$ & $+15''{,}43$ & $0{,}998\,8839$ & $-0''{,}49$ & $23^{0}\,27'\,59''{,}48$\\
17 & $204\quad20\quad21{,}54$ & $+15{,}51$ & $0{,}995\,3968$ & $+0{,}79$ & $59{,}26$\\
27 & $214\quad16\quad52{,}21$ & $+15{,}60$ & $0{,}992\,8340$ & $-0{,}15$ & $59{,}06$\\
\hline
\end{tabular}
\end{center}

We shall set up the calculation as though the orbit were still entirely unknown. We must therefore transfer the parallax to the positions of the Earth, since we cannot remove it from the positions of Juno. First we reduce the observed positions themselves from the equator to the ecliptic, using the apparent obliquity. This gives the following values.
\begin{center}
\small\setlength{\tabcolsep}{6pt}\renewcommand{\arraystretch}{1.2}
\begin{tabular}{|r|c|c|}
\hline
 & \shortstack{App. longitude\\of Juno} & \shortstack{App. latitude\\of Juno}\\
\hline
Oct. 5 & $354^{0}\,44'\,54''{,}28$ & $-4^{0}\,59'\,31''{,}59$\\
17 & $352\quad34\quad44{,}51$ & $-6\quad21\quad56{,}25$\\
27 & $351\quad34\quad51{,}57$ & $-7\quad17\quad52{,}70$\\
\hline
\end{tabular}
\end{center}

\Needspace{14\baselineskip}
We now combine this calculation with the determination of the longitude and latitude of the zenith at the observing site for the three observations. The right ascension agrees with Juno's right ascension (because the observations were made on the meridian itself), and the declination equals the altitude of the pole $=51^{0}\,28'\,39''$.
Thus we obtain
\begin{center}
\small\setlength{\tabcolsep}{6pt}\renewcommand{\arraystretch}{1.2}
\begin{tabular}{|r|c|c|}
\hline
 & Longitude of the zenith & Latitude\\
\hline
Oct. 5 & $24^{0}\,29'$ & $46^{0}\,53'$\\
17 & $23\quad25$ & $47\quad24$\\
27 & $23\quad1$ & $47\quad36$\\
\hline
\end{tabular}
\end{center}
\noindent\SourceLeaf{202}Following the rules given in Art. 72, we now determine the fictitious positions of the Earth in the plane of the ecliptic itself, at which the celestial body would have appeared just as it did at the actual positions of observation. Taking the mean solar parallax $=8''{,}6$ gives
\begin{center}
\small\setlength{\tabcolsep}{4pt}\renewcommand{\arraystretch}{1.2}
\begin{tabular}{|r|c|c|c|}
\hline
 & Longitude reduction & Distance reduction & Time reduction\\
\hline
Oct. 5 & $-22''{,}39$ & $+0{,}000\,3856$ & $-0^{\mathrm{s}}{,}19$\\
17 & $-27{,}21$ & $+0{,}000\,2329$ & $-0{,}12$\\
27 & $-35{,}82$ & $+0{,}000\,2085$ & $-0{,}12$\\
\hline
\end{tabular}
\end{center}
\noindent The time reduction is included only to show that it is entirely imperceptible.

Next all the longitudes, both of the planet and of the Earth, must be reduced to the mean vernal equinox for a chosen epoch, for which we adopt the beginning of 1805. After subtracting the nutation, we therefore still have to add the precession, which for the three observations is $11''{,}87$, $10''{,}23$, $8''{,}86$ respectively. We must thus add $-3''{,}56$ for the first observation, $-5''{,}28$ for the second, and $-6''{,}74$ for the third.

Finally, Juno's longitudes and latitudes must be corrected for the aberration of the fixed stars. By the known rules we find that $19''{,}12$, $17''{,}11$, $14''{,}82$ are to be subtracted from the respective longitudes, and $0''{,}53$, $1''{,}18$, $1''{,}75$ are to be added to the respective latitudes. This addition reduces the absolute values, since southern latitudes are treated as negative.
